\section{Stationary height slopes}\label{sec:stationary}

The volume estimates of the preceding sections detect which linear maps
can occur on the height quotient.  Our first source of linear maps is
reduced cohomology.  We apply it to an invariant probability on a compact
space of normalized quasi-isometries.  This gives an averaged sublinear
estimate for one fixed measure.  A coupling of the forward and inverse
maps then shows that its slope lies in the finite group $\mathcal A$ of
Lemma~\ref{lem:support-symmetries}.  Removing the invariant measure is a
separate task, beginning in Section~\ref{sec:profiles}.

The use of invariant measures on spaces of quasi-isometries and their
associated cocycles follows the measured-coupling approach developed by
Shalom \cite[Sections~2.1--2.2]{Shalom2004}. The continuous pair space and
the transfer of the same slope to its inverse are constructed below.

Throughout this section $G$ has the hypotheses of
Lemma~\ref{lem:cartan-model}, and $W\ne0$.  We write $e$ for its identity.
The homomorphism $\pi:G\to W$ is Lipschitz after a Euclidean norm on $W$
is fixed: its differential at the identity has bounded operator norm,
and left translation transports that bound to every point.

\subsection{Lipschitz pairs and their compact sections}

We first replace coarse maps by continuous maps without losing
uniformity.  This lets us use compact-open limits and continuous group
actions.

\begin{lemma}[Lipschitz representatives]\label{lem:lipschitz-pairs}
For each $K\ge1$ and $C\ge0$, there are constants $B<\infty$ and
$\Lambda\ge1$, depending only on $G,K,C$, with the following property.
Every $(K,C)$ self quasi-isometry of $G$ is within distance $B$ of a map
$F$ belonging to a pair $(F,\bar F)$ such that both maps are
$\Lambda$-Lipschitz, both satisfy
\[
 d(J(g),J(h))\ge \Lambda^{-1}d(g,h)-\Lambda
 \qquad(J=F,\bar F),
\]
and
\[
 d(\bar F(F(g)),g)\le\Lambda,\qquad
 d(F(\bar F(g)),g)\le\Lambda
 \qquad(g\in G).
\]
\end{lemma}

The proof is given in Appendix~\ref{app:lipschitz-representatives}.

Fix $\Lambda$ as in Lemma~\ref{lem:lipschitz-pairs}.  Let $\Omega$ be
the space of all pairs satisfying its displayed bounds, with the
compact-open topology on both maps.  More generally, everything below
applies to any nonempty closed subspace invariant under the two
translations that we now define.  Target translation and domain
translation act by
\begin{align*}
 L_y(F,\bar F)&=(L_y\circ F,\bar F\circ L_{y^{-1}}),\\
 (F,\bar F)\cdot x&=(F\circ L_x,L_{x^{-1}}\circ\bar F),
\end{align*}
where the $L$'s inside the pairs are ordinary left translations on $G$.
These actions commute and preserve every bound defining $\Omega$.
Define the normalized sections
\[
 X=\{(F,\bar F)\in\Omega:F(e)=e\},\qquad
 Y=\{(F,\bar F)\in\Omega:\bar F(e)=e\}.
\]

\begin{lemma}[The two sections]\label{lem:pair-sections}
The spaces $X$ and $Y$ are compact metrizable.  There are
homeomorphisms
\[
 G\times X\longrightarrow\Omega,\quad (z,\xi)\longmapsto L_z\xi,
 \qquad
 Y\times G\longrightarrow\Omega,\quad (\eta,x)\longmapsto\eta\cdot x.
\]
The formulas
\[
 T_x\xi=L_{F_\xi(x)^{-1}}(\xi\cdot x),\qquad
 U_y\eta=(L_{y^{-1}}\eta)\cdot\bar F_\eta(y)
\]
define continuous right actions on $X$ and $Y$, respectively.
In particular, $T_yT_x=T_{xy}$ and $U_zU_y=U_{yz}$.
\end{lemma}

\begin{proof}
On $X$, the composition bound gives
$d(\bar F(e),e)\le\Lambda$, since $F(e)=e$.  On $Y$ it similarly
gives $d(F(e),e)\le\Lambda$.  Equi-Lipschitz bounds now give pointwise
precompactness on every compact subset of $G$.  Arzel\`a--Ascoli and a
countable compact exhaustion yield compactness; all the inequalities
and composition bounds are closed under compact-open convergence.
The same exhaustion metrizes the topology.

The inverse of the first coordinate map is
\[
 \omega\longmapsto
 \bigl(F_\omega(e),L_{F_\omega(e)^{-1}}\omega\bigr),
\]
and that of the second is
\[
 \omega\longmapsto
 \bigl(\omega\cdot\bar F_\omega(e),\bar F_\omega(e)^{-1}\bigr).
\]
They are continuous.  Continuity of evaluation and composition follows
directly from compact-open convergence and the common Lipschitz bound.
The action identities follow either from these coordinates or by
substituting the two normalization formulas.  For example,
\[
 F_{T_x\xi}(z)=F_\xi(x)^{-1}F_\xi(xz),\qquad
 \bar F_{U_y\eta}(z)=\bar F_\eta(y)^{-1}\bar F_\eta(yz).
\]
These identities also verify the stated right-action convention.
\end{proof}

Here and below an invariant probability means invariant under every
element of $G$.  Such probabilities exist on any nonempty compact
invariant section: $G$ is solvable and therefore amenable, so applying
an invariant mean to an orbit of each continuous function gives an
invariant positive functional; its representing probability is
invariant.  Extreme points of the compact convex set of invariant
probabilities are ergodic.  We will also use the ergodic decomposition
of a fixed invariant probability, which applies because the action is
continuous, $G$ is second countable, and the section is compact
metrizable; see \cite[Theorem~5.2]{GreschonigSchmidt2000}, specialized
to the invariant case. Inversion converts our right action to the
left-action convention in that theorem.

\subsection{Reduced cohomology and one forward comparison}

For an invariant probability $\mu$ on $X$, the height evaluation is the
$L^2(X,\mu;W)$-valued cocycle
\[
 b(x)(\xi)=\pi F_\xi(x),\qquad
 (\rho(x)f)(\xi)=f(T_x\xi).
\]
The convention in Lemma~\ref{lem:pair-sections} gives
$\rho(x)\rho(y)=\rho(xy)$, and
\begin{equation}\label{eq:height-cocycle}
 b(xy)=b(x)+\rho(x)b(y).
\end{equation}

\begin{lemma}[A slope for a fixed ergodic measure]
\label{lem:cohomological-slope}
Let $\mu$ be an invariant ergodic probability on $X$.  There is a
continuous homomorphism $l:G\to W$ and functions
$\phi_j\in L^2(X,\mu;W)$ such that, for every compact $Q\subset G$,
\begin{equation}\label{eq:compact-coboundaries}
 \sup_{x\in Q}
 \bigl\|b(x)+\phi_j\circ T_x-\phi_j-l(x)\bigr\|_{L^2(\mu)}
 \longrightarrow0.
\end{equation}
The homomorphism $l$ kills $N$ and factors through $V$ on $D$.
For each fixed $0<R<\infty$,
\begin{equation}\label{eq:stationary-sublinear}
 \lim_{r\to\infty}\frac1r\sup_{|x|_G\le Rr}
 \|\pi F(x)-l(x)\|_{L^2(\mu)}=0.
\end{equation}
These assertions concern the fixed measure $\mu$; no rate uniform
over invariant measures is asserted.
\end{lemma}

\begin{proof}
The representation $\rho$ is unitary because $\mu$ is invariant.
It is strongly continuous: this follows first on continuous functions
from the continuous action on compact $X$, and then on $L^2$ by
density and unitarity.  The jointly continuous evaluation
$(x,\xi)\mapsto\pi F_\xi(x)$ is uniformly continuous on compact
subsets of $G\times X$, so $b$ is a continuous $L^2$ cocycle.

Real-triangulable Lie groups have property $\mathcal{WAP}_{\mathrm t}$,
and hence property $H_{\mathrm t}$, by
\cite[Corollary~1.9]{CornulierTessera2020}.  In the form needed here,
every continuous unitary representation without invariant vectors has
zero first reduced cohomology.  Reduced cohomology uses the closure
of the coboundaries for uniform norm convergence on compact subsets
of the group.  The Hilbert representations here are weakly almost
periodic, so the cited theorem applies.

Ergodicity identifies the invariant vectors in $L^2(X,\mu;W)$ with
the constant functions.  Orthogonal projection of
\eqref{eq:height-cocycle} onto this subspace gives the continuous
homomorphism
\[
 l(x)=\int_X\pi F_\xi(x)\,d\mu(\xi).
\]
Apply the vanishing theorem to the cocycle $b-l$ in the orthogonal
complement.  It gives \eqref{eq:compact-coboundaries}.  One may
complexify the real Hilbert space and take real parts if using the
complex convention for unitary representations.

For completeness, fix a compact symmetric generating neighborhood
$Q$ such that every $x$ is a product of at most $a|x|_G+a$ elements
of $Q$.  A fixed-radius metric ball has this property by subdividing
a geodesic.  Set
\[
 c_j(x)=b(x)+\rho(x)\phi_j-\phi_j-l(x).
\]
Given $\epsilon>0$, fix $j$ with
$\sup_{x\in Q}\|c_j(x)\|_2\le\epsilon$.
The cocycle identity and unitarity imply
\[
 \|b(x)-l(x)\|_2
 \le (a|x|_G+a)\epsilon+2\|\phi_j\|_2.
\]
Keep this $j$ fixed, let $r\to\infty$ for $|x|_G\le Rr$, and then
let $\epsilon\downarrow0$.  This proves
\eqref{eq:stationary-sublinear}.  It is a supremum of $L^2$ norms,
which will be sufficient for finite meshes of paths.

A continuous homomorphism to the additive group $W$ kills the
derived Lie algebra.  Since $\mathfrak n\subset[\mathfrak g,\mathfrak g]$
and $N$ is connected, it kills $N$; on $D$ it factors through the
abelianization $V$.
\end{proof}

We now use the first volume comparison to show that this homomorphism
detects exactly the exponential height directions.  The argument does
not require ergodicity, provided the averaged estimate is already
available.

\begin{lemma}[Volume comparison for an averaged slope]
\label{lem:stationary-volume}
Let $\sigma$ be a probability on $X$, and let $l:G\to W$ be a
continuous homomorphism killing $N$.  Suppose
\eqref{eq:stationary-sublinear} holds with $\sigma$ in place of
$\mu$.  Regard $l|_D$ as a linear map on $V$.  Then every nonempty
compact convex $E_V\subset V$ satisfies
\begin{equation}\label{eq:forward-slope-volume}
 P(\pi E_V)\le P(lE_V).
\end{equation}
Consequently $l=A\pi$ for an invertible linear map $A:W\to W$.
The same assertion holds for the inverse maps on $Y$.
\end{lemma}

\begin{proof}
Fix $v_0\in E_V$ and a pivot $d_0\in D$ of intrinsic length $O(r)$
whose abelianization is $rv_0$.  Put $E=\pi E_V$.  Apply
Lemma~\ref{lem:support-switch-paths} with the set $E_V\subset V$:
although $l$ factors through $V$ on $D$, we do not yet know that it
factors through $W$.  For each $n\in B_r(E)$, that lemma supplies a
measurably chosen path from $d_0$ to $nd_0$ of length $O(r)$.
Its abelianized $D$-legs stay in $rE_V$, and its switches have length
$o(r)$.  Since $l$ kills $N$ and is Lipschitz, the path's $l$-heights
lie in $r(lE_V+o(1))$.  All its points have distance $O(r)$ from $e$,
with common bounds for the sampled endpoints.

Sample $n$ with normalized Haar measure on $B_r(E)$, independently
of $\xi\sim\sigma$.  Fix $\tau>0$ and sample each path at an
arclength mesh of spacing at most $\tau r$, including its endpoints.
There are at most $M_\tau$ sampled points, independent of $r,n$.
At any of them, say $g_r(n)$, the assumed estimate gives
\[
 \int_X|\pi F_\xi(g_r(n))-l(g_r(n))|^2\,d\sigma(\xi)=o(r^2),
\]
uniformly in $n$.  Fubini, a finite union bound, and Chebyshev's
inequality show that, for every fixed $\eta>0$, the fraction of pairs
$(\xi,n)$ on which a sampled error exceeds $\eta r$ tends to zero.
Hence for all large $r$ there is one $\xi$ for which at least half
the endpoint volume has no such error.

For these endpoints the entire image path lies at heights within
$(c\tau+\eta+o(1))r$ of $r lE_V$.  Indeed the image of each mesh
interval has diameter at most the common Lipschitz constant times
$\tau r$.  The initial image point is the common $F_\xi(d_0)$;
subtracting its height only translates the allowed height set.

At their common source $D$-position, the good endpoints have a
separated subset of cardinality at least
$c_0 e^{rP(E)}$, by Lemma~\ref{lem:fixed-sheet-packing}.
Choose the separation large enough that the quasi-isometry lower
bound keeps their images separated by a fixed positive amount.
Their image paths have length $O(r)$.  Lemma~\ref{lem:path-packing}
and translation invariance of $P$ give
\[
 P(E)\le P\bigl(lE_V+\overline B_W(0,c\tau+\eta)\bigr)
\]
after taking $r\to\infty$.  First let $\eta\downarrow0$, then
$\tau\downarrow0$.  Continuity of $P$ for Hausdorff distance proves
\eqref{eq:forward-slope-volume}.  This argument only used a fixed
finite mesh before the limit; it did not require an $L^2$ bound on
the pointwise supremum over a path.

Apply \eqref{eq:forward-slope-volume} to $[0,v]\subset V$.
Since $p$ is a norm,
\[
 p(lv)\ge p(\pi v),\qquad \ker l\subseteq\ker\pi.
\]
Writing $d=\dim W$, the inclusion forces $\rank l\ge d$, while
the target of $l$ has dimension $d$.  Thus $l$ and $\pi$ have the
same kernel on $V$, and $l=A\pi$ there for an invertible $A$.
Both homomorphisms kill $N$ and $G=ND$, so this equality holds on $G$.
All the reasoning applies to the reversed pairs on $Y$ as well.
\end{proof}

\subsection{Transferring the same slope to the inverse}

An independently obtained slope for $\bar F$ would not yet show
that $A$ preserves $P$.  We need its slope to be exactly $A^{-1}$.
The commuting actions on $\Omega$ give a finite measure on the
inverse section, and averaging over a ball makes the transfer functions
well-defined there.

\begin{lemma}[The inverse coupling]\label{lem:inverse-coupling}
Let $\mu$ be an invariant probability on $X$ and let $A\in\GL(W)$.
Suppose there are $\phi_j\in L^2(X,\mu;W)$ satisfying
\eqref{eq:compact-coboundaries} with $l=A\pi$.
There is a finite nonzero $U$-invariant measure $\nu$ on $Y$ and
functions $\theta_j\in L^2(Y,\nu;W)$ such that, for every compact
$Q\subset G$,
\begin{equation}\label{eq:inverse-coboundaries}
 \sup_{y\in Q}
 \bigl\|\pi\bar F(y)+\theta_j\circ U_y-\theta_j-A^{-1}\pi y
 \bigr\|_{L^2(\nu)}\longrightarrow0.
\end{equation}
In particular, for every fixed $0<R<\infty$,
\begin{equation}\label{eq:inverse-stationary-sublinear}
 \lim_{r\to\infty}\frac1r\sup_{|y|_G\le Rr}
 \|\pi\bar F(y)-A^{-1}\pi y\|_{L^2(\nu)}=0.
\end{equation}
\end{lemma}

\begin{proof}
\emph{The invariant measure on the inverse section.}
Let $m$ be Haar measure on $G$.  In the coordinates $\Omega=G\times X$
of Lemma~\ref{lem:pair-sections}, put $M=m\times\mu$.
This is a nonzero Radon measure and is target-invariant.  Domain
translation has the formula
\[
 (z,\xi)\cdot x=(zF_\xi(x),T_x\xi).
\]
For each fixed $\xi$, the change in $z$ is right translation,
which preserves $m$ by unimodularity.  Invariance of $\mu$ therefore
makes $M$ domain-invariant as well.

In the other coordinates $\Omega=Y\times G$, domain translation is
$(\eta,x)\mapsto(\eta,xg)$.  For a Borel set $E\subset Y$, the
measure $B\mapsto M(E\times B)$ is right-invariant and locally
finite: for compact $B$ it is bounded by $M(Y\times B)<\infty$.
Haar uniqueness gives $M(E\times B)=\nu(E)m(B)$, and countable
additivity in $E$ supplies a Borel measure $\nu$ with
\begin{equation}\label{eq:two-section-measure}
 M=\nu\times m\quad\hbox{on }Y\times G.
\end{equation}
Here a product with compact factors is compact in $\Omega$, by the
coordinate homeomorphism, so the local finiteness used above follows
from the Radon property of $M$.

Choose a compact ball $B_0\subset G$ with $0<m(B_0)<\infty$.
Then
\[
 \nu(Y)m(B_0)=M(Y\times B_0)<\infty.
\]
If $\nu(Y)=0$, a countable compact exhaustion of $G$ and
\eqref{eq:two-section-measure} would imply $M=0$.
Thus $0<\nu(Y)<\infty$.  In particular, $\nu$ is a measure obtained
from product coordinates; we are not restricting $M$ to the
possibly $M$-null section $Y$.

For $y\in G$, set $q_y(\eta)=\bar F_\eta(y)$.  The identity
\begin{equation}\label{eq:coupling-target-action}
 L_{y^{-1}}(\eta\cdot x)
   =(U_y\eta)\cdot(q_y(\eta)^{-1}x)
\end{equation}
shows that $\nu$ is $U$-invariant.  More explicitly, integrate a
test function of $\eta$ times a compactly supported function of $x$
against $M$.  In the transformed integral, the fiber change
$x\mapsto q_y(\eta)^{-1}x$ preserves left Haar measure.
Target invariance of $M$ then gives invariance of the $\eta$-integral.

\emph{Transfer functions averaged over the fibers.}
Choose measurable representatives of the $\phi_j$ and lift them by
\[
 \Phi_j(L_z\xi)=\pi z+\phi_j(\xi).
\]
They belong to $L^2$ on every compact subset of $\Omega$ and satisfy
$\Phi_j(L_z\omega)=\pi z+\Phi_j(\omega)$.
Their domain increments have defect
\[
 E_j(\omega,t)=\Phi_j(\omega\cdot t)-\Phi_j(\omega)-A\pi t.
\]
If $\omega=L_z\xi$, this is exactly
$c_j(t)(\xi)=b(t)(\xi)+\phi_j(T_t\xi)-\phi_j(\xi)-A\pi t$.
Consequently, for compact $K_\Omega\subset\Omega$ and
$K_G\subset G$,
\begin{equation}\label{eq:local-increment-defect}
 \sup_{t\in K_G}\int_{K_\Omega}|E_j(\omega,t)|^2\,dM(\omega)
 \longrightarrow0.
\end{equation}
Indeed $K_\Omega$ lies in $K\times X$ for a compact $K\subset G$
in the first coordinates; the integral is at most
$m(K)\|c_j(t)\|_2^2$.

We now average over whole balls so that no trace of an $L^2$ function
on a null section is required.  Fix a compact ball $B\subset G$
of positive Haar measure, and define
\begin{align*}
 D_j(\eta,x)&=\Phi_j(\eta\cdot x)-A\pi x,\\
 \Psi_j(\eta)&=\frac1{m(B)}\int_B D_j(\eta,x)\,dm(x).
\end{align*}
Product coordinates and Fubini define this average for $\nu$-almost
every $\eta$.  Jensen's inequality and local $L^2$ integrability of
$\Phi_j$ show that $\Psi_j\in L^2(Y,\nu;W)$; define it arbitrarily
on the remaining null set.

Equation~\eqref{eq:coupling-target-action}, target equivariance of
$\Phi_j$, and additivity of $\pi$ give
\begin{equation}\label{eq:inverse-transfer-identity}
 \Psi_j(U_y\eta)-\Psi_j(\eta)
 =-\pi y+A\pi q_y(\eta)+R_j(\eta,y),
\end{equation}
where $R_j(\eta,y)$ is the difference of the averages of
$D_j(\eta,\cdot)$ on $q_y(\eta)B$ and on $B$.
For each fixed $y$, this identity holds almost everywhere:
$U_y$ preserves $\nu$, and the changes in the fiber variable preserve
Haar measure.  These facts also show that changing the measurable
representatives of the transfers changes none of their $L^2$ classes
or averages.

\emph{Comparing two fiber averages.}
We claim that $R_j$ tends to zero in $L^2(\nu)$ uniformly for $y$
in any fixed compact $Q$.  The increment estimate
\eqref{eq:local-increment-defect} will control, in $L^2(\nu)$, the
difference between averages over two bounded regions of each fiber.  Since
$\bar F_\eta(e)=e$ on $Y$ and the
maps have a common Lipschitz constant, all $q_y(\eta)$ with
$y\in Q$, $\eta\in Y$ lie in one compact set $Q_1\subset G$.
Put $C_1=Q_1B$.  Jensen's inequality applied to the two averages
gives, for every $y\in Q$,
\begin{align}
 \|R_j(\cdot,y)\|_2^2
 &\le \frac1{m(B)^2}
 \int_Y\int_{C_1}\int_B
 |D_j(\eta,x_1)-D_j(\eta,x_2)|^2\,
 dm(x_2)\,dm(x_1)\,d\nu(\eta).
 \label{eq:whole-ball-bound}
\end{align}
The enlargement from $q_y(\eta)B$ to $C_1$ is legitimate because
its averaging density is bounded by $1/m(B)$, even though its center
depends on $\eta$.  Thus the right side is independent of $y$.

Inside this integral one has
\[
 D_j(\eta,x_1)-D_j(\eta,x_2)
   =E_j(\eta\cdot x_2,x_2^{-1}x_1).
\]
Make the substitution $t=x_2^{-1}x_1$, which preserves left Haar
measure in $x_1$.  Enlarge the resulting $t$-domain to the compact
$B^{-1}C_1$.  The pairs $\eta\cdot x_2$ range over the compact
image of $Y\times B$, and their measure is $M$ by
\eqref{eq:two-section-measure}.  Hence the right side of
\eqref{eq:whole-ball-bound} is at most
\[
 \frac1{m(B)^2}\int_{B^{-1}C_1}
 \int_{Y\cdot B}|E_j(\omega,t)|^2\,dM(\omega)\,dm(t),
\]
which tends to zero by \eqref{eq:local-increment-defect}.
We have proved
\[
 \sup_{y\in Q}\|R_j(\cdot,y)\|_{L^2(\nu)}\longrightarrow0.
\]
The supremum here is over the individual $L^2$ norms.  No common
pointwise exceptional set for all $y$ is being asserted or used.

Set $\theta_j=-A^{-1}\Psi_j$ and rearrange
\eqref{eq:inverse-transfer-identity}.  This gives
\[
 \pi\bar F_\eta(y)+\theta_j(U_y\eta)-\theta_j(\eta)
 =A^{-1}\pi y-A^{-1}R_j(\eta,y),
\]
proving \eqref{eq:inverse-coboundaries} with precisely the inverse
of the given $A$.  The inverse height evaluations form a continuous
cocycle for the unitary action on $L^2(Y,\nu;W)$, by the formulas
in Lemma~\ref{lem:pair-sections}.  Subdividing a geodesic into bounded
steps, exactly as in the proof of
Lemma~\ref{lem:cohomological-slope}, proves
\eqref{eq:inverse-stationary-sublinear}.  Neither ergodicity of $\nu$
nor a second application of reduced cohomology is needed.
\end{proof}

\subsection{Finite stationary labels}

\begin{proposition}\label{prop:stationary-slopes}
For every invariant ergodic probability $\mu$ on $X$, there is an
$A_\mu\in\mathcal A$ such that
\eqref{eq:compact-coboundaries} holds with $l=A_\mu\pi$ and
\[
 \lim_{r\to\infty}\frac1r\sup_{|x|_G\le Rr}
 \|\pi F(x)-A_\mu\pi x\|_{L^2(\mu)}=0
 \qquad(0<R<\infty).
\]
The coupled finite invariant measure on $Y$ satisfies the analogous
estimate with $A_\mu^{-1}$.
\end{proposition}

\begin{proof}
Lemma~\ref{lem:cohomological-slope} gives $l$ and its compact
coboundary approximations.  Lemma~\ref{lem:stationary-volume} gives
$l=A\pi$ with $A$ invertible.  The inverse coupling in
Lemma~\ref{lem:inverse-coupling} supplies a finite nonzero invariant
measure $\nu$ and the inverse averaged estimate with the same
$A^{-1}$.  Normalize $\nu$ to a probability and apply
Lemma~\ref{lem:stationary-volume} to the inverse maps.

Every compact convex $E\subset W$ has a compact convex lift to $V$,
by a fixed linear section of $\pi:V\to W$.  The forward comparison
therefore gives $P(E)\le P(AE)$, and the inverse comparison applied
to $AE$ gives $P(AE)\le P(E)$.  Thus $P(AE)=P(E)$ for every
nonempty compact convex $E$, which is precisely $A\in\mathcal A$.
\end{proof}

The next form allows an occupation measure to have several ergodic
components.  It is the only stationary conclusion needed for the
microstep argument in the next section.

\begin{corollary}[Stationary increments]\label{cor:stationary-increments}
Fix an invariant probability $\mu$ on $X$, without assuming ergodicity.
For each $v\in V$, represented in the fixed degree-one complement
in $\mathfrak d$, one has
\begin{equation}\label{eq:finite-stationary-increments}
 \int_X\dist\left(
   \frac{\pi F_\xi(\exp_D(Lv))}{L},
   \{A\pi v:A\in\mathcal A\}\right)\,d\mu(\xi)
 \longrightarrow0\qquad(L\to\infty).
\end{equation}
For each fixed $L>0$, the integrand is a bounded continuous function
on $X$.  For $L\ge1$ these functions have a common bound depending
only on $v,G,\Lambda$.  The convergence is for the fixed $\mu$.
\end{corollary}

\begin{proof}
Write $\mu=\int\mu_s\,d\kappa(s)$ for its ergodic decomposition
\cite[Theorem~5.2]{GreschonigSchmidt2000}.
On almost every component Proposition~\ref{prop:stationary-slopes}
gives $A_s\in\mathcal A$.  Since
$|\exp_D(Lv)|_G\le L\|v\|$ up to the fixed norm convention, its
estimate implies convergence of the displayed integral with $\mu_s$
in place of $\mu$.

The normalized height evaluations are uniformly bounded for $L\ge1$
by the Lipschitz bounds for $F$ and $\pi$ and by $F(e)=e$.
The comparison set is finite and fixed.  Dominated convergence
through the ergodic decomposition proves
\eqref{eq:finite-stationary-increments}.  No measurable choice of
$A_s$ is required: the integrand already takes distance to the full
finite set.  Continuity for fixed $L$ follows from evaluation in
the compact-open topology.
\end{proof}

When applying Corollary~\ref{cor:stationary-increments} to limits of
occupation measures, one must first fix a limiting invariant measure.
One may then pass to the geometric limit for each fixed microstep
length $L$, using continuity of its integrand, and only afterwards
let $L\to\infty$.  This order uses exactly the fixed-measure
convergence proved here.
